\bta{2019}

\begin{enumerate}
	%\renewcommand{\labelenumi}{\arabic{enumi}.}
	% A(\Alph) a(\alph) I(\Roman) i(\roman) 1(\arabic)
	%设定全局标号series=example	%引用全局变量resume=example
	%[topsep=-0.3em,parsep=-0.3em,itemsep=-0.3em,partopsep=-0.3em]
	%可使用leftmargin调整列表环境左边的空白长度 [leftmargin=0em]
	\item
设$ 1 $维线性谐振子的初始状态为\(\psi(x, 0)=A\left[3 \psi_{0}(x)+4 \psi_{1}(x)\right]\)，其中
	\(\psi_{0}(x), \psi_{1}(x)\)分别为体系的归一化的基态和第一激发态的波函数。
	\begin{enumerate}
		%\renewcommand{\labelenumi}{\arabic{enumi}.}
		% A(\Alph) a(\alph) I(\Roman) i(\roman) 1(\arabic)
		%设定全局标号series=example	%引用全局变量resume=example
		%[topsep=-0.3em,parsep=-0.3em,itemsep=-0.3em,partopsep=-0.3em]
		%可使用leftmargin调整列表环境左边的空白长度 [leftmargin=0em]
		\item
求归一化常数$ A $。
	
	\item 
	计算\(\psi(x, t)\)、\(|\psi(x, t)|^{2}\)。
	
	\item 
	求粒子能量的可取值，各可取值出现的概率，及能量的平均值。
		\end{enumerate}

\banswer{
	\begin{enumerate}
		\item 由 $|A|^2(3^2+4^2)=1$，可取
		\[
			\boxed{A=\frac15}.
		\]

		\item 令 $E_0=\hbar\omega/2$、$E_1=3\hbar\omega/2$，则
		\[
			\boxed{\psi(x,t)=\frac15\left[
			3\psi_0(x)e^{-iE_0t/\hbar}
			+4\psi_1(x)e^{-iE_1t/\hbar}\right]}.
		\]
		取实的定态波函数后，
		\[
			\boxed{|\psi(x,t)|^2=\frac1{25}\left[
			9\psi_0^2(x)+16\psi_1^2(x)
			+24\psi_0(x)\psi_1(x)\cos(\omega t)\right]}.
		\]

		\item 能量 $E_0$、$E_1$ 的概率分别为 $9/25$、$16/25$，故
		\[
			\boxed{\langle H\rangle=\frac9{25}\frac{\hbar\omega}{2}
			+\frac{16}{25}\frac{3\hbar\omega}{2}
			=\frac{57}{50}\hbar\omega}.
		\]
	\end{enumerate}
}


\item 
设体系由$ 3 $个无相互作用粒子组成，$ 3 $个粒子分别处于正交归一单粒子态\(\psi_a\),\(\psi_b\),\(\psi_c\)。
\begin{enumerate}
	%\renewcommand{\labelenumi}{\arabic{enumi}.}
	% A(\Alph) a(\alph) I(\Roman) i(\roman) 1(\arabic)
	%设定全局标号series=example	%引用全局变量resume=example
	%[topsep=-0.3em,parsep=-0.3em,itemsep=-0.3em,partopsep=-0.3em]
	%可使用leftmargin调整列表环境左边的空白长度 [leftmargin=0em]
	\item
若$ 3 $个粒子为自旋为$ 0 $的非全同玻色子，求出其波函数。
	
	\item 
	若$ 3 $个粒子为自旋为$ 0 $的全同玻色子，求出其波函数。
	
	\item 
	若$ 3 $个粒子为自旋$ 1/2 $的全同费米子，自旋处于\(|\uparrow\rangle\)，求出其波函数。
\end{enumerate}


\banswer{
	以下用 $\psi_\alpha(j)$ 表示粒子 $j$ 处于单粒子态 $\psi_\alpha$。
	\begin{enumerate}
		\item 对可区分的三个粒子，若已指定粒子 $1,2,3$ 分别处于 $a,b,c$，则
		\[
			\boxed{\Psi_{\rm d}(1,2,3)=\psi_a(1)\psi_b(2)\psi_c(3)}.
		\]

		\item 对全同自旋零玻色子，波函数必须完全对称。因三个轨道互异且正交，
		\[
			\boxed{\Psi_B=\frac1{\sqrt{3!}}\sum_{P\in S_3}
			\psi_{P(a)}(1)\psi_{P(b)}(2)\psi_{P(c)}(3)}.
		\]

		\item 三个费米子的自旋均为 $|\uparrow\rangle$，自旋部分
		$|\uparrow\uparrow\uparrow\rangle$ 完全对称，故空间部分必须完全反对称：
		\[
			\boxed{\Psi_F=\frac1{\sqrt{3!}}
			\begin{vmatrix}
			\psi_a(1)&\psi_b(1)&\psi_c(1)\\
			\psi_a(2)&\psi_b(2)&\psi_c(2)\\
			\psi_a(3)&\psi_b(3)&\psi_c(3)
			\end{vmatrix}
			|\uparrow\uparrow\uparrow\rangle}.
		\]
	\end{enumerate}
}


\item 
设粒子的轨道角动量量子数\(l=1\)（取\(\hbar=1\)）。
	\begin{enumerate}
		%\renewcommand{\labelenumi}{\arabic{enumi}.}
		% A(\Alph) a(\alph) I(\Roman) i(\roman) 1(\arabic)
		%设定全局标号series=example	%引用全局变量resume=example
		%[topsep=-0.3em,parsep=-0.3em,itemsep=-0.3em,partopsep=-0.3em]
		%可使用leftmargin调整列表环境左边的空白长度 [leftmargin=0em]
		\item
写出\((\hat{L}^{2},\hat{L}_z)\)表象中\(\hat{L_y}\)的矩阵表示。
	
	\item 
	证明\(\hat{L}_{y}^{3}=\hat{L}_{y}\)。
	
	\item 
	证明算符\(e^{-i\lambda\hat{L}_{y}}=1-i\hat{L}_{y}sin\lambda-(1-cos\lambda)\hat{L}_{y}^2\)。
	
	\item 
	求出\(e^{-i\lambda\hat{L}_{y}}\)在\((\hat{L}^{2},\hat{L}_y)\)表象中的矩阵表示。
	\end{enumerate}
\banswer{
	在按 $m_z=1,0,-1$ 排列的基底中（$\hbar=1$）：
	\begin{enumerate}
		\item
		\[
			\boxed{L_y=\frac1{\sqrt2}
			\begin{pmatrix}0&-i&0\\ i&0&-i\\0&i&0\end{pmatrix}}.
		\]

		\item $L_y$ 的本征值为 $1,0,-1$，故多项式
		$x^3-x$ 在其每个本征值上均为零。由谱定理
		\[
			\boxed{L_y^3=L_y}.
		\]
		也可直接将上一矩阵相乘验证。

		\item 利用 $L_y^{2k+1}=L_y$（$k\ge0$）及
		$L_y^{2k}=L_y^2$（$k\ge1$），指数级数化为
		\begin{align*}
			e^{-i\lambda L_y}
			&=I+\left(\sum_{k=0}^{\infty}
			\frac{(-i\lambda)^{2k+1}}{(2k+1)!}\right)L_y
			+\left(\sum_{k=1}^{\infty}
			\frac{(-i\lambda)^{2k}}{(2k)!}\right)L_y^2\\
			&=\boxed{I-iL_y\sin\lambda-(1-\cos\lambda)L_y^2}.
		\end{align*}

		\item 在 $(L^2,L_y)$ 表象中 $L_y=\operatorname{diag}(1,0,-1)$，所以
		\[
			\boxed{e^{-i\lambda L_y}
			=\begin{pmatrix}
			e^{-i\lambda}&0&0\\
			0&1&0\\
			0&0&e^{i\lambda}
			\end{pmatrix}}.
		\]
	\end{enumerate}
}


\newpage	
	\item 
	一个质量为m的粒子处在一维无限深方势阱\(\left(-\frac{L}{2}<x<\frac{L}{2}\right)\)中，其中哈密顿量记为\(H_0\)。
	\begin{enumerate}
		%\renewcommand{\labelenumi}{\arabic{enumi}.}
		% A(\Alph) a(\alph) I(\Roman) i(\roman) 1(\arabic)
		%设定全局标号series=example	%引用全局变量resume=example
		%[topsep=-0.3em,parsep=-0.3em,itemsep=-0.3em,partopsep=-0.3em]
		%可使用leftmargin调整列表环境左边的空白长度 [leftmargin=0em]
		\item
求\(H_0\)的本征值和对应的波函数\(\psi(x)\)。
	
	\item 
	记受到微扰\(H^{\prime}(x)=\lambda \delta(x)(\lambda>0)\)，求基态能量和基态波函数的一级修正。
	
	\item 
	严格求\(H=H_{0}+H'\)下的能级方程，并考虑\(\lambda \ll \frac{2 \hbar^{2}}{m L}\)的极限并与上述一级微扰论的基态能量比较。
\end{enumerate}	

\banswer{
	\begin{enumerate}
		\item 令 $\epsilon=\pi^2\hbar^2/(2mL^2)$。可取
		\[
			E_n^{(0)}=n^2\epsilon,\qquad
			\phi_n(x)=\sqrt{\frac2L}
			\begin{cases}
			\cos(n\pi x/L),&n\ {\rm odd},\\
			\sin(n\pi x/L),&n\ {\rm even},
			\end{cases}
		\]
		其中 $n=1,2,\ldots$ 且 $|x|<L/2$。

		\item 基态能量的一级修正为
		\[
			\boxed{E_1^{(1)}=\lambda|\phi_1(0)|^2=\frac{2\lambda}{L}}.
		\]
		只有偶宇称态在原点非零，故基态波函数的一级修正为
		\[
			\boxed{|\phi_1^{(1)}\rangle
			=-\frac{m\lambda L}{\pi^2\hbar^2}
			\sum_{j=1}^{\infty}\frac{|\,2j+1\,\rangle}{j(j+1)}}.
		\]
		因此 $|\psi_{\rm g}\rangle=|1\rangle+|\phi_1^{(1)}\rangle+O(\lambda^2)$；
		一级精度下不需要额外归一化修正。

		\item 奇宇称态（$n$ 为偶数）在 $x=0$ 有节点，$\delta$ 势对它们不起作用，
		故 $E_n=n^2\epsilon$ 严格成立。偶宇称态在 $0<x<L/2$ 可写为
		$\psi=C\sin[k(L/2-x)]$，它在 $x=L/2$ 处自动为零。把薛定谔方程在
		$(-\eta,\eta)$ 上积分后令 $\eta\to0$，得导数跃变条件
		\[
			\psi'(0^+)-\psi'(0^-)=\frac{2m\lambda}{\hbar^2}\psi(0).
		\]
		代入 $\psi$ 得 $-2k\cos\frac{kL}{2}=\frac{2m\lambda}{\hbar^2}\sin\frac{kL}{2}$，
		即严格能级方程
		\[
			\boxed{k\cot\frac{kL}{2}=-\frac{m\lambda}{\hbar^2}},\qquad
			E=\frac{\hbar^2k^2}{2m}.
		\]
		这正是通常的 $\delta$ 势束缚条件的对称化形式：偶宇称态等价于
		$\tan\frac{kL}{2}=-\frac{\hbar^2k}{m\lambda}$，奇宇称态则完全不满足它。
		在弱耦合极限（$\lambda\ll2\hbar^2/(mL)$，即无量纲小量
		$\beta=\frac{m\lambda L}{2\hbar^2}\ll1$）令第 $j$ 个偶宇称态
		$k=\frac{(2j-1)\pi}{L}+\delta k_j$，即 (1) 中 $n=2j-1$，基态为 $j=1$。
		该支解位于 $\frac{(2j-1)\pi}{2}<\frac{kL}{2}<j\pi$，正是由无扰第
		$2j-1$ 个能级连续变形而来的那一支。由
		$\cot\left(j\pi-\frac{\pi}{2}+x\right)=-\tan x$ 得
		\[
			\left[\frac{(2j-1)\pi}{L}+\delta k_j\right]
			\tan\frac{\delta k_jL}{2}=\frac{m\lambda}{\hbar^2}.
		\]
		取小角近似 $\tan x\simeq x$（要求
		$\frac{\delta k_jL}{2}=\frac{2\beta}{(2j-1)\pi}\ll1$，即
		$\lambda\ll\frac{(2j-1)\pi\hbar^2}{mL}$，与题设条件同量级），并在系数中
		只保留零阶项，得到
		\[
			\delta k_j=\frac{2m\lambda}{(2j-1)\pi\hbar^2}+O(\lambda^2).
		\]
		代回 $E=\hbar^2k^2/2m$ 即得
		\[
			E_j=\frac{(2j-1)^2\pi^2\hbar^2}{2mL^2}+\frac{2\lambda}{L}+O(\lambda^2).
		\]
		一阶能量移动 $2\lambda/L$ 与 $j$ 无关，因为各偶宇称态在原点取值都是
		$\sqrt{2/L}$。取基态 $j=1$：
		\[
			\boxed{E_{\rm g}=\frac{\pi^2\hbar^2}{2mL^2}+\frac{2\lambda}{L}+O(\lambda^2)},
		\]
		与一级微扰论结果 $E_1^{(1)}=\lambda|\phi_1(0)|^2=\frac{2\lambda}{L}$ 完全一致。
		继续展开可得二阶项 $-\frac{2m\lambda^2}{(2j-1)^2\pi^2\hbar^2}$
		（$j=1$ 时为 $-\frac{2m\lambda^2}{\pi^2\hbar^2}$），与二级微扰论相符。
	\end{enumerate}
}


\item 
由$ 3 $个电子组成的体系，设其自旋哈密顿量为\((\hbar=1)\)
$ \hat{H}=a(\hat{S_1^x}\hat{S_2^x}+\hat{S_1^y}\hat{S_2^y}+\hat{S_2^y}\hat{S_3^y}+\hat{S_2^x}\hat{S_3^x}+\hat{S_1^x}\hat{S_3^x}+\hat{S_1^y}\hat{S_3^y})-b\hat{S_T^z} $
其中\(\hat{S_i^\alpha}\)代表第$ i $个电子自旋算符的$\alpha $
分量，$\hat{S_T^z}$代表体系总自旋算符$\vec{S_T}=\vec{S_1}+\vec{S_2} + \vec{S_3}$的$ z $分量。
\begin{enumerate}
	%\renewcommand{\labelenumi}{\arabic{enumi}.}
	% A(\Alph) a(\alph) I(\Roman) i(\roman) 1(\arabic)
	%设定全局标号series=example	%引用全局变量resume=example
	%[topsep=-0.3em,parsep=-0.3em,itemsep=-0.3em,partopsep=-0.3em]
	%可使用leftmargin调整列表环境左边的空白长度 [leftmargin=0em]
	\item
求出总自旋平方算符\(\hat{S_T^2}\)的最大和最小本征值
	
	\item 
	计算以下两个对易式：\(\left[\hat{S_{T}^{2}}, \hat{H}\right]\)、\(\left[\hat{S_{T}^{z}}, \hat{H}\right]\)
	
	\item 
	求出\(\hat{H}\)的所有本征值。
\end{enumerate}

	
		\banswer{
			把相互作用项写成总自旋算符：
			\[
				\sum_{i<j}(S_i^xS_j^x+S_i^yS_j^y)
				=\frac12\left(S_T^2-(S_T^z)^2-\frac32\right).
			\]
			因此
			\[
				\boxed{H=\frac a2\left(S_T^2-(S_T^z)^2-\frac32\right)
				-bS_T^z}.
			\]
			\begin{enumerate}
				\item 三个自旋 $1/2$ 的总自旋可为 $S=3/2$ 或 $S=1/2$，所以
				\[
					\boxed{(S_T^2)_{\max}=\frac{15}{4},\qquad
					(S_T^2)_{\min}=\frac34}.
				\]

				\item 从上面的 $H(S_T^2,S_T^z)$ 形式立即得到
				\[
					\boxed{[S_T^2,H]=0,\qquad[S_T^z,H]=0}.
				\]

				\item 一般能量为
				\[
					E_{S,M}=\frac a2\left[S(S+1)-M^2-\frac32\right]-bM.
				\]
				对 $S=3/2$ 四重态：
				\[
					\boxed{
					E_{\frac32,\frac32}=-\frac{3b}{2},\quad
					E_{\frac32,\frac12}=a-\frac b2,\quad
					E_{\frac32,-\frac12}=a+\frac b2,\quad
					E_{\frac32,-\frac32}=\frac{3b}{2}}.
				\]
				对 $S=1/2$：
				\[
					\boxed{E_{\frac12,\frac12}=-\frac a2-\frac b2,\qquad
					E_{\frac12,-\frac12}=-\frac a2+\frac b2}.
				\]
				$S=1/2$ 表示在三个自旋的直积空间中出现两次，因此这两个能量各有二重简并；连同四重态共得到全部八个态。
			\end{enumerate}
		}
	
	
	
\end{enumerate}

